\section{Lições de Outras Áreas da Saúde}

A odontologia não opera em um vácuo epistemológico e não precisa, portanto, trilhar caminhos já explorados por outras disciplinas médicas e das engenharias. O amadurecimento dos gêmeos digitais exige uma visão interdisciplinar.

Na cardiologia, a validação de gêmeos digitais do miocárdio, utilizados para predição de fluxos hemodinâmicos e ablação de arritmias, alcançou um nível de maturidade superior (FDA-approved em alguns contextos). A principal lição herdada da medicina cardiovascular é a obrigatoriedade da \textbf{validação contínua}: cada predição \textit{in silico} deve ser realimentada pelo desfecho clínico, constituindo um ciclo de aprendizado supervisionado contínuo (\textit{closed-loop learning}) \cite{katsoulakis2024digital}. 

A ortopedia, por sua vez, oferta um severo aviso (\textit{cautionary tale}). Nos primórdios do planejamento virtual de artroplastias de quadril e joelho, a utilização de modelos ósseos linear-elásticos que ignoravam o comportamento viscoelástico complexo dos tecidos moles circunjacentes resultou em predições cinemáticas implausíveis. Quando transferidas para a cirurgia (via guias estereolitográficos primitivos), essas falhas de modelagem levaram a taxas de luxação e soltura asséptica inaceitáveis. Na odontologia, subestimar a complexidade do ligamento periodontal na simulação de ortodontia com alinhadores ou do tecido gengival em prótese resultará no mesmo fracasso biomecânico \cite{li2025impacts}.

Finalmente, no seu berço aeroespacial (NASA), instituiu-se que a utilidade de um gêmeo digital decai exponencialmente se ele não for alimentado por telemetria (\textit{real-time sensor feedback}). A transposição desse conceito exige que o gêmeo odontológico seja dinâmico, atualizado a cada nova consulta, e não uma imagem congelada do dia do escaneamento \cite{menon2026digital}. A falta de biossensores intraorais de baixo custo, entretanto, ainda torna esse sincronismo contínuo um horizonte distante para a rotina odontológica.
